\chapter{Metaorders, Latent Liquidity and Cross-Impact}\label{ch:mx:metaorders-latent-liquidity-and-cross-impact}

Why should impact grow like the square root of size, whatever the stock, the decade or the market? The answer most models now give is that the book a
trader sees is a small part of the liquidity that exists: the rest is latent, waiting to be revealed at prices a large order has not yet reached. This
chapter builds that latent book, recovers the square root and the linear regime below it, checks when a metaorder is priced fairly, and turns to the
impact that trading one stock has on another.

\section{Why impact is concave}

Chapter~11 measured a concave impact and chapter~12 a transient one. Three explanations of the concavity are in use. Splitting: large metaorders are cut
into children whose impacts decay between them. Information: Farmer, Gerig, Lillo and Waelbroeck (2013) derived the shape of impact from competition among
execution services and a condition of fair pricing. Liquidity: Toth and co-authors (2011, chapter~11) argued that the average supply and demand vanish at
the price and grow linearly away from it. The third has become a model.

\section{Latent liquidity and the locally linear book}

\begin{definition}[Latent liquidity, latent order book]\label{def:mx:metaorders-latent-liquidity-and-cross-impact:latent}
\emph{Latent liquidity}\index{latent liquidity} is the intention to trade at a price that has not been placed in the book, because its owner waits
for the price to come near. The \emph{latent order book}\index{latent order book} is the density of these intentions by price, of which the visible book
is the small part closest to the price.
\end{definition}

\begin{definition}[Locally linear order book]\label{def:mx:metaorders-latent-liquidity-and-cross-impact:linear}
A \emph{locally linear order book}\index{locally linear order book} has a latent density that vanishes at the price and grows linearly away from it,
$\varphi(x)=L\,|x-p|$, with $L$ the liquidity per unit of price squared.
\end{definition}

Donier, Bonart, Mastromatteo and Bouchaud (2015) derived it from the reaction--diffusion of latent orders: buyers' and sellers' intentions diffuse as
their valuations change, and annihilate (trade) where they meet. The signed density $\varphi$ (asks positive) then obeys a diffusion equation whose
stationary solution is linear. A buy metaorder at rate $m$ removes asks at the price, and the price, the zero of $\varphi$, moves up. If the metaorder is
fast compared with diffusion, it simply eats the linear book: moving the price by $I$ consumes $\int_0^IL\,x\,dx=LI^2/2$, so $I=\sqrt{2Q/L}$. If it is
slow, the book refills by diffusion as it trades, and impact is linear in size. Bucci, Benzaquen, Lillo and Bouchaud (2019) documented the crossover.

\texttt{firm.crossimpact.llob} integrates the diffusion on a grid and removes the metaorder's shares from the asks nearest the price at each step.

\omcode{code/firm/crossimpact/firm_crossimpact.py}{33}{56}{The latent order book: the signed density diffuses, the metaorder removes the nearest asks, and the price is the zero of the density.}\label{lst:mx:metaorders-latent-liquidity-and-cross-impact:llob}

With $L=1$, $D=1$ and a ten-second execution, the \omterm{def:mx:the-empirics-of-price-impact:peak}{peak impact} is 0.05 ticks for $Q=0.3$, 0.18 for 1, 0.52 for 3, 1.74 for 10, 4.9 for 30, 12.2 for 100, 23.3 for 300 and
44.0 for 1\,000 (\cref{fig:mx:metaorders-latent-liquidity-and-cross-impact:scan}). The local exponent is 1.0 up to $Q=10$, then 0.94, 0.76, 0.59 and 0.53: linear
below $LDT=10$, a square root above, where the impact approaches $\sqrt{2Q/L}$ (24.5 for 300, 44.7 for 1\,000).

\begin{omfigure}
\begin{tikzpicture}
\begin{axis}[width=9cm,height=5.4cm,xmode=log,ymode=log,xlabel={\small metaorder size $Q$},ylabel={\small peak impact (ticks)},
  tick label style={font=\footnotesize},grid=major,grid style={gray!25},legend columns=3,legend style={font=\footnotesize,at={(0.5,-0.3)},
  anchor=north,draw=none,fill=none,/tikz/every even column/.append style={column sep=8pt}},table/col sep=comma]
\addplot[omDef,very thick,mark=*] table[x=q,y=peak]{figdata/microstructure/13-metaorders-latent-liquidity-and-cross-impact/scan.csv};
\addplot[gray,dashed,domain=0.3:30] {0.176*x};
\addplot[omThm,dashed,domain=30:1000] {sqrt(2*x)};
\legend{latent book,linear,$\sqrt{2Q/L}$}
\end{axis}
\end{tikzpicture}

\omcaption{\omterm{def:mx:the-empirics-of-price-impact:peak}{Peak impact} of a ten-second metaorder on a locally linear latent book ($L=1$, $D=1$): linear in size when the book refills faster than the
order eats it, a square root when it does not. Data: \texttt{mx\_cross.llob\_scan}.}\label{fig:mx:metaorders-latent-liquidity-and-cross-impact:scan}
\end{omfigure}

\section{Fair pricing and the end of a metaorder}

\begin{definition}[Fair pricing condition]\label{def:mx:metaorders-latent-liquidity-and-cross-impact:fair}
The \emph{fair pricing condition}\index{fair pricing condition} (Farmer, Gerig, Lillo and Waelbroeck, 2013) states that the average price paid for a
metaorder equals the price after it is completed: whoever filled it was not paid for being picked off, and the metaorder's owner paid for what its trading
revealed.
\end{definition}

If impact grows like the square root of the time traded, the average impact paid over the execution is two thirds of the peak, so fair pricing predicts
that the price settles at two thirds of its peak once the metaorder ends. On the latent book, a metaorder of 10 reaches 1.74 ticks and pays 1.16 on average,
0.66 of the peak (\cref{fig:mx:metaorders-latent-liquidity-and-cross-impact:path}); after the end the price falls through the average paid 1.9 seconds
later, and keeps falling, to 0.16 of the peak after 100 seconds. For a metaorder of 300 the average is again 0.66 of the peak, crossed 2.3 seconds after
the end, with 0.27 left after 100 seconds. The model without memory loss makes impact transient, so the price is fair only for an instant; in the full
model, finite memory of latent orders leaves a permanent part.

\begin{omfigure}
\begin{tikzpicture}
\begin{axis}[width=9cm,height=5cm,xlabel={\small seconds},ylabel={\small impact (ticks)},tick label style={font=\footnotesize},xmin=0,xmax=110,
  ymin=0,ymax=2,grid=major,grid style={gray!25},table/col sep=comma]
\fill[gray!12] (axis cs:0,0) rectangle (axis cs:10,2);
\addplot[omDef,very thick] table[x=t,y=p]{figdata/microstructure/13-metaorders-latent-liquidity-and-cross-impact/path.csv};
\addplot[omThm,dashed,domain=0:110] {1.155};
\node[omThm,font=\footnotesize,anchor=south west] at (axis cs:30,1.16) {average price paid};
\end{axis}
\end{tikzpicture}

\omcaption{A metaorder of 10 on the latent book: the impact during the execution (shaded) and after it, against the average price paid. The price
crosses the average two seconds after the end and keeps falling. Data: \texttt{mx\_cross.fair}.}\label{fig:mx:metaorders-latent-liquidity-and-cross-impact:path}
\end{omfigure}

\section{Cross-impact}

\begin{definition}[Cross-impact, cross-impact matrix]\label{def:mx:metaorders-latent-liquidity-and-cross-impact:cross}
\emph{Cross-impact}\index{cross-impact} is the effect of trading one asset on the price of another. The \emph{cross-impact matrix}\index{cross-impact matrix}
$\Lambda$ gives each asset's price change per unit of each asset's order flow, $r_t=\Lambda q_t+\text{noise}$.
\end{definition}

Pasquariello and Vega (2015) found in US stocks that one industry's or stock's daily order imbalance moves other industries' and stocks' returns,
persistently and often negatively. Benzaquen, Mastromatteo, Eisler and Bouchaud (2017) built a multivariate propagator in which liquidity shares the
sectorial structure of correlations, and found that trades mediate a significant part of the covariance of returns.

Two simulated stocks with order flows correlated 0.6 (thousands of shares per interval) and a true matrix $\Lambda=\begin{pmatrix}1.0&0.4\\0.4&0.8\end{pmatrix}$
ticks per thousand shares, with unit noise, give over 2\,000 intervals the estimate $\begin{pmatrix}0.96&0.42\\0.42&0.81\end{pmatrix}$. Beyond its own flow,
the other stock's flow explains 4.5\% of the first stock's return variance and 5.3\% of the second's: correlated flows already carry most of the
common move, and \omterm{def:mx:metaorders-latent-liquidity-and-cross-impact:cross}{cross-impact} adds the rest.

\section{Arbitrage-free cross-impact}

Schneider and Lillo (2019) extended the no-dynamic-arbitrage argument to several assets: with bounded \omterm{def:mx:transient-impact-and-propagator-models:propagator}{decay kernels}, \omterm{def:mx:metaorders-latent-liquidity-and-cross-impact:cross}{cross-impact} must be linear and odd in
trading intensity and symmetric, the impact of $i$ on $j$ equal to that of $j$ on $i$; otherwise a trader could cycle through the assets at a profit.
An estimated matrix is not exactly symmetric; \texttt{symmetric\_psd} projects it on the nearest symmetric positive semi-definite matrix before any
optimiser uses it.

What does \omterm{def:mx:metaorders-latent-liquidity-and-cross-impact:cross}{cross-impact} change for execution? Selling 10 (thousand shares) of each stock over one unit of time with transient impact decaying at rate 3,
the cheapest joint schedule costs 52.1 (\cref{fig:mx:metaorders-latent-liquidity-and-cross-impact:costs}). An own-impact model, which ignores the
cross terms, estimates the same schedule at 36.0, 31\% too low. For a pair trade, selling 10 of the first and buying 10 of the second, the joint optimum costs
20.0 and the own-impact estimate is 36.0, 80\% too high: each leg's \omterm{def:mx:metaorders-latent-liquidity-and-cross-impact:cross}{cross-impact} helps the other. When \omterm{def:mx:metaorders-latent-liquidity-and-cross-impact:cross}{cross-impact} decays like own impact, the optimal
schedules are the same with or without it; what the model gets wrong is the cost. Schedules differ when the legs are traded one after the other, as a
desk that thinks the legs independent may do: then the joint schedule saves 19\% of the sequential cost for the sale and 48\% for the pair trade. When
\omterm{def:mx:metaorders-latent-liquidity-and-cross-impact:cross}{cross-impact} decays more slowly than own impact (0.3 against 3), the joint optimum also changes shape, and saves a further 3.7\% on the pair trade.

\omcode{code/firm/crossimpact/firm_crossimpact.py}{92}{111}{Joint liquidation under transient own- and cross-impact, or asset by asset ignoring the cross terms.}\label{lst:mx:metaorders-latent-liquidity-and-cross-impact:liq}

\begin{omfigure}
\begin{tikzpicture}
\begin{axis}[width=8.5cm,height=5cm,ybar,bar width=12pt,xtick={0,1},xticklabels={sell both,pair trade},ylabel={\small cost},
  tick label style={font=\footnotesize},ymin=0,ymax=75,grid=major,grid style={gray!25},enlarge x limits=0.4,
  legend columns=3,legend style={font=\footnotesize,at={(0.5,-0.2)},anchor=north,draw=none,fill=none,
  /tikz/every even column/.append style={column sep=8pt}},table/col sep=comma]
\addplot[omDef,fill=omDef!40] table[x=k,y=joint]{figdata/microstructure/13-metaorders-latent-liquidity-and-cross-impact/costs.csv};
\addplot[omProp,fill=omProp!40] table[x=k,y=own]{figdata/microstructure/13-metaorders-latent-liquidity-and-cross-impact/costs.csv};
\addplot[omThm,fill=omThm!40] table[x=k,y=sequential]{figdata/microstructure/13-metaorders-latent-liquidity-and-cross-impact/costs.csv};
\legend{joint optimum,own-impact estimate,legs one after the other}
\end{axis}
\end{tikzpicture}

\omcaption{Selling two stocks, or trading them as a pair, under \omterm{def:mx:metaorders-latent-liquidity-and-cross-impact:cross}{cross-impact}: the cost of the joint optimum, what an own-impact model estimates for it,
and the cost of trading the legs in turn. Data: \texttt{mx\_cross.cross\_study}.}\label{fig:mx:metaorders-latent-liquidity-and-cross-impact:costs}
\end{omfigure}

\section{Tutorial: selling the stock, moving its neighbour}\label{tut:mx:metaorders-latent-liquidity-and-cross-impact:tut}

\begin{tutorial}
\textbf{Goal.} Recover the square root from a latent book, check fair pricing, and estimate and use \omterm{def:mx:metaorders-latent-liquidity-and-cross-impact:cross}{cross-impact}.
\textbf{End state:} \cref{fig:mx:metaorders-latent-liquidity-and-cross-impact:scan,fig:mx:metaorders-latent-liquidity-and-cross-impact:path,fig:mx:metaorders-latent-liquidity-and-cross-impact:costs}.
\begin{tutsteps}
\item \textbf{Latent book.} \texttt{firm\_crossimpact.llob(q, horizon, L, D, dt, width, dx, after)}; \texttt{mx\_cross.llob\_scan()} over eight sizes.
\item \textbf{Fair pricing.} \texttt{fair\_pricing(times, price, horizon)} and \texttt{mx\_cross.fair(q)}.
\item \textbf{\omterm{def:mx:metaorders-latent-liquidity-and-cross-impact:cross}{Cross-impact}.} \texttt{two\_assets()}, \texttt{estimate(r, q)}, \texttt{symmetric\_psd}, \texttt{explained\_by\_other}.
\item \textbf{Liquidation.} \texttt{liquidation(Lambda, rho, times, X, joint, rho\_cross)} for a sale and a pair trade; draw with \texttt{fig\_cross.py}.
\end{tutsteps}
\textbf{What to change next.} Add cancellation and deposition of latent orders (finite memory) and look for a permanent impact; estimate \omterm{def:mx:metaorders-latent-liquidity-and-cross-impact:cross}{cross-impact} on
three assets and decompose it on the eigenvectors of the flows' covariance.
\end{tutorial}

\section{Build: latent liquidity and cross-impact}\label{bld:mx:metaorders-latent-liquidity-and-cross-impact:crossimpact}

\begin{build}
\bfield{Purpose} The square-root law's mechanism and the multi-asset impact that portfolio executions (chapters 15 and 22) and Book~11's multi-asset
market making need.
\bfield{Interface} \texttt{llob(q, horizon, L, D, dt, width, dx, after)}, \texttt{fair\_pricing(times, price, horizon)}, \texttt{estimate(r, q)},
\texttt{symmetric\_psd(Lambda)}, \texttt{explained\_by\_other(r, q)}, \texttt{liquidation(Lambda, rho, times, X, joint, rho\_cross)}.
\bfield{Rules} Explicit diffusion (stable for $D\,dt/dx^2\le\tfrac12$) with fixed linear far ends; the price is the interpolated zero of the density;
\omterm{def:mx:metaorders-latent-liquidity-and-cross-impact:cross}{cross-impact} by least squares of returns on flows; symmetric projection by eigenvalue clipping.
\bfield{Acceptance tests} \texttt{code/firm/crossimpact/tests/}: a fast metaorder moves the price by about $\sqrt{2Q/L}$; a planted matrix recovered;
the projection symmetric and positive; the joint liquidation no dearer than asset by asset.
\bfield{Stretch} Latent orders with deposition and cancellation; the eigenliquidity decomposition; nonlinear \omterm{def:mx:metaorders-latent-liquidity-and-cross-impact:cross}{cross-impact}.
\end{build}

\begin{omsources}
\item J.~D.~Farmer, A.~Gerig, F.~Lillo and H.~Waelbroeck, ``How efficiency shapes market impact'', \emph{Quantitative Finance} 13(11), 2013.
\item J.~Donier, J.~Bonart, I.~Mastromatteo and J.-P.~Bouchaud, ``A fully consistent, minimal model for non-linear market impact'', \emph{Quantitative
Finance} 15(7), 2015.
\item P.~Pasquariello and C.~Vega, ``Strategic cross-trading in the U.S. stock market'', \emph{Review of Finance} 19(1), 2015.
\item M.~Benzaquen, I.~Mastromatteo, Z.~Eisler and J.-P.~Bouchaud, ``Dissecting cross-impact on stock markets: an empirical analysis'', \emph{Journal of
Statistical Mechanics}, 2017.
\item M.~Schneider and F.~Lillo, ``Cross-impact and no-dynamic-arbitrage'', \emph{Quantitative Finance} 19(1), 2019.
\item F.~Bucci, M.~Benzaquen, F.~Lillo and J.-P.~Bouchaud, ``Crossover from linear to square-root market impact'', \emph{Physical Review Letters} 122, 2019.
\end{omsources}

\section{Exercises}

\begin{exercise}[$\star$]\label{exo:mx:metaorders-latent-liquidity-and-cross-impact:1}
A locally linear latent book has $L=2$ (thousand shares per tick squared). How far does a fast metaorder of 50 thousand shares move the price?
\end{exercise}

\begin{exercise}[$\star$]\label{exo:mx:metaorders-latent-liquidity-and-cross-impact:2}
Why is impact linear in size when the metaorder is slow compared with the latent book's diffusion?
\end{exercise}

\begin{exercise}[$\star$]\label{exo:mx:metaorders-latent-liquidity-and-cross-impact:3}
Show that if impact grows like $\sqrt t$ during an execution at a constant rate, the average impact paid is two thirds of the peak.
\end{exercise}

\begin{exercise}[$\star\star$]\label{exo:mx:metaorders-latent-liquidity-and-cross-impact:4}
The estimated \omterm{def:mx:metaorders-latent-liquidity-and-cross-impact:cross}{cross-impact matrix} is $\begin{pmatrix}1.0&0.5\\-0.1&-0.2\end{pmatrix}$. What is wrong with it, and what does the projection do?
\end{exercise}

\begin{exercise}[$\star\star$]\label{exo:mx:metaorders-latent-liquidity-and-cross-impact:5}
Why does an own-impact model underestimate the cost of selling two correlated stocks and overestimate that of a pair trade?
\end{exercise}

\begin{exercise}[$\star\star$]\label{exo:mx:metaorders-latent-liquidity-and-cross-impact:6}
Why is the joint optimal schedule the same with and without \omterm{def:mx:metaorders-latent-liquidity-and-cross-impact:cross}{cross-impact} when both decay at the same rate?
\end{exercise}

\begin{exercise}[$\star\star\star$]\label{exo:mx:metaorders-latent-liquidity-and-cross-impact:7}
\emph{Coding.} Rerun the latent-book scan with a horizon of 100 seconds instead of 10. Where does the crossover move, and why?
\end{exercise}

\begin{exercise}[$\star\star\star$]\label{exo:mx:metaorders-latent-liquidity-and-cross-impact:8}
\emph{Find the flaw.} ``The price fell back below our average fill two seconds after we finished, so the market makers who filled us lost money.''
\end{exercise}

\section{Problem: Selling the Stock, Moving Its Neighbour}

\begin{problem}[{Weekend problem --- selling the stock, moving its neighbour}]\label{pb:mx:metaorders-latent-liquidity-and-cross-impact:1}
A portfolio manager must sell two correlated stocks and later run a pair trade in them. Explain the impact and plan the executions.

\medskip\noindent\textbf{Part I --- \omterm{def:mx:metaorders-latent-liquidity-and-cross-impact:latent}{Latent liquidity}.}
\begin{enumerate}
\item Give the three explanations of concave impact.
\item Define the latent and the locally linear book.
\item Derive $\sqrt{2Q/L}$ for a fast metaorder.
\item Give the simulated \omterm{def:mx:the-empirics-of-price-impact:peak}{peak impacts} and local exponents, and where the crossover is.
\end{enumerate}

\medskip\noindent\textbf{Part II --- Fair pricing.}
\begin{enumerate}[resume]
\item State the \omterm{def:mx:metaorders-latent-liquidity-and-cross-impact:fair}{fair pricing condition}.
\item Solve exercise~3.
\item When does the simulated price cross the average paid, and what happens next?
\item Why is impact transient in this model, and what would make part of it permanent?
\end{enumerate}

\medskip\noindent\textbf{Part III --- \omterm{def:mx:metaorders-latent-liquidity-and-cross-impact:cross}{Cross-impact}.}
\begin{enumerate}[resume]
\item Define the \omterm{def:mx:metaorders-latent-liquidity-and-cross-impact:cross}{cross-impact matrix}.
\item What did Pasquariello and Vega, and Benzaquen and co-authors, find?
\item Give the estimated matrix and the explained shares.
\item State Schneider and Lillo's conditions and why they are needed.
\end{enumerate}

\medskip\noindent\textbf{Part IV --- Execution and the verdict.}
\begin{enumerate}[resume]
\item Give the joint cost and the own-impact estimate for the sale.
\item And for the pair trade.
\item What does trading the legs one after the other cost?
\item When does \omterm{def:mx:metaorders-latent-liquidity-and-cross-impact:cross}{cross-impact} change the optimal schedule itself?
\item State the \emph{named result}: the share of each asset's price move explained by the other asset's order flow, and the cost saved by a
cross-impact-aware liquidation.
\item How should the manager execute the sale?
\item And the pair trade?
\item In one sentence: why does the neighbour matter?
\end{enumerate}
\end{problem}

\section{Interview questions}

\begin{interviewq}[$\star$ \iqroles{researcher}]\label{iq:mx:metaorders-latent-liquidity-and-cross-impact:1}
Explain the square-root law with the \omterm{def:mx:metaorders-latent-liquidity-and-cross-impact:latent}{latent order book} in three sentences.
\end{interviewq}

\begin{interviewq}[$\star\star$ \iqroles{researcher}]\label{iq:mx:metaorders-latent-liquidity-and-cross-impact:2}
What is the \omterm{def:mx:metaorders-latent-liquidity-and-cross-impact:fair}{fair pricing condition}, and what does it predict for the price after a metaorder?
\end{interviewq}

\begin{interviewq}[$\star\star$ \iqroles{trader}]\label{iq:mx:metaorders-latent-liquidity-and-cross-impact:3}
You sell a large position in a stock. Which other positions of yours does it hurt, and how would you know?
\end{interviewq}

\begin{interviewq}[$\star\star$ \iqroles{researcher}]\label{iq:mx:metaorders-latent-liquidity-and-cross-impact:4}
How would you estimate a \omterm{def:mx:metaorders-latent-liquidity-and-cross-impact:cross}{cross-impact matrix}, and what constraints would you impose?
\end{interviewq}

\begin{interviewq}[$\star\star$ \iqroles{developer}]\label{iq:mx:metaorders-latent-liquidity-and-cross-impact:5}
Simulate a diffusing density on a grid with a moving boundary. What limits the time step?
\end{interviewq}

\begin{interviewq}[$\star\star\star$ \iqroles{trader, researcher}]\label{iq:mx:metaorders-latent-liquidity-and-cross-impact:6}
Why might a pair trade cost less than the sum of its legs, and how would you exploit that in execution?
\end{interviewq}
